
\setcounter{footnote}{0}
\SGARefTarget{sga3:xviii:chapter:xviii}\chapter*{Exposé XVIII\\
Weil's Theorem on the Construction of a Group from a Rational Law}
\addcontentsline{toc}{chapter}{Exposé XVIII. Weil's Theorem on the
Construction of a Group from a Rational Law}
\markboth{EXPOSÉ XVIII. WEIL'S THEOREM}
  {EXPOSÉ XVIII. WEIL'S THEOREM}

\begin{center}
  \textit{by Michael Artin}
\end{center}
\begingroup
\renewcommand{\thefootnote}{0}
\footnotetext{Version of May 8, 2009.}
\endgroup

\SGARefTarget{sga3:xviii:section:0}\section*{0. Introduction}
\addcontentsline{toc}{section}{0. Introduction}

This Exposé is devoted to the well-known theorem of Weil~[1], which
constructs a group variety from a birational law. It seems that the
generalization of this result to the case where the base is the
spectrum of a discrete valuation ring was already known to several
people; see, for example,~[2]. Here we shall prove the theorem for a
flat group of finite presentation over an arbitrary base prescheme
\(S\). Since the hypotheses must be formulated somewhat more
carefully, the statement is not in the form given by Weil; but when
\(S\) is the spectrum of a field, the form given here is seen to be
essentially equivalent to Weil's. We use the following suggestion of
Grothendieck: given a ``group germ'' \(X\) over a prescheme \(S\), one
should construct the group as the quotient of \(X\times_S X\) by a
suitable equivalence relation.

\SGARefTarget{sga3:xviii:section:1}\section*{1. ``Recollections'' on
Rational Maps}
\addcontentsline{toc}{section}{1. ``Recollections'' on Rational Maps}

Let \(X/S\) be a relative prescheme and let \(U\subseteq X\) be an
open subset. We say that \(U\) is \emph{schematically dense in \(X\)
relative to \(S\)}, or that \(U\hookrightarrow X\) is
\emph{relatively schematically dense}, if, for every base change
\(S'\to S\), the open subset \(U\times_S S'\) is schematically dense
in \(X\times_S S'\). For the definition and properties of this
notion, we refer to Exposé IX, \S4.\begingroup
\renewcommand{\thefootnote}{*}\footnote{See also EGA
IV\textsubscript{3}, 11.9 and 11.10 (especially 11.10.8), where the
term ``universally schematically dense relative to \(S\)'' is used.}
\addtocounter{footnote}{-1}
\endgroup

\SGARefTarget{sga3:xviii:prop:1.1}\paragraph{\SGARefLink{sga3:xviii:prop:1.1}{Proposition 1.1}.}
\textit{(i) A finite intersection, as well as a union of a nonempty
family, of open subsets schematically dense relative to \(S\), is
schematically dense relative to \(S\).}

\textit{(ii) If \(U\subseteq X\) is schematically dense relative to
\(S\), and if \(S\to T\) is a morphism, then
\(U_T\subseteq X_T\) is schematically dense relative to \(T\).}

\textit{(iii) Let \(U\subseteq V\subseteq X\) be open immersions. For
\(U\) to be schematically dense in \(X\) relative to \(S\), it is
necessary and sufficient that it be so in \(V\), and that \(V\) be so
in \(X\).}

\textit{\SGAInlineRefTarget{sga3:xviii:prop:1.1:item:iv}(iv) If \(U\subseteq X\) and \(V\subseteq Y\) are relatively
schematically dense, then \(U\times_S V\) is schematically dense in
\(X\times_S Y\) relative to \(S\).}

In Exposé IX one finds the following criterion:

\SGARefTarget{sga3:xviii:prop:1.2}\paragraph{\SGARefLink{sga3:xviii:prop:1.2}{Proposition 1.2}.}
\textit{Let \(X/S\) be locally of finite presentation and flat, and
let \(U\) be an open subset of \(X\). If, for every \(s\in S\), the
fiber \(U_s\) is schematically dense in \(X_s\), then \(U\) is
schematically dense in \(X\) relative to \(S\).}

In particular:

\SGARefTarget{sga3:xviii:cor:1.3}\paragraph{\SGARefLink{sga3:xviii:cor:1.3}{Corollary 1.3}.}
\textit{Suppose that \(X/S\) is locally of finite presentation and
flat, and that every fiber \(X_s\) is ``without embedded
components.'' Then \(U\) is schematically dense in \(X\) relative to
\(S\) if and only if, for every \(s\in S\), \(U_s\) is dense in
\(X_s\) in the topological sense.}

One easily deduces from the definition the following proposition.

\SGARefTarget{sga3:xviii:prop:1.4}\paragraph{\SGARefLink{sga3:xviii:prop:1.4}{Proposition 1.4}.}
\textit{Let \(U\subseteq X\) be schematically dense relative to \(S\),
and let \(f,g:X\to Y\) be two morphisms of functors over \(S\), where
\(Y\) satisfies one of the following conditions:}

\begin{enumerate}
\item[\SGAInlineRefTarget{sga3:xviii:prop:1.4:item:i}(i)] \textit{\(Y\) is a prescheme over \(S\), and every fiber
  \(Y_s\) is separated.}
\item[\SGAInlineRefTarget{sga3:xviii:prop:1.4:item:ii}(ii)] \textit{\(Y\) is a presheaf\footnote{Editor's note:
  specify for which topology; a priori, the fpqc topology.} over
  \(S\), such that the diagonal morphism
  \(Y\to Y\times_S Y\) is representable by a closed
  immersion.\footnote{Editor's note: Let us recall here the
  definition of a closed immersion. A morphism of \(S\)-functors
  \(F\to G\) is \emph{relatively representable} if, for every
  \(S\)-morphism \(T\to G\), where \(T\) is an \(S\)-prescheme, the
  \(T\)-functor \(F_T=F\times_G T\) is representable by a
  \(T\)-prescheme. The French source prints \(F\times_S T\) here;
  the fiber product over \(G\) is the one dictated by the preceding
  morphism \(T\to G\). A morphism of \(S\)-functors \(F\to G\) is an open
  (respectively, closed) immersion if it is relatively representable
  and if, for every \(S\)-morphism \(T\to G\), where \(T\) is an
  \(S\)-prescheme, the \(T\)-morphism \(F_T\to T\) is an open
  (respectively, closed) immersion.}}
\end{enumerate}

\textit{Then, if \(f=g\) on \(U\), one has \(f=g\).}

\textit{Proof.}
In case \SGARefLink{sga3:xviii:prop:1.4:item:i}{(i)}, the standard argument following IX, \SGARefLink{sga3:ix:definition:4:1}{4.1} shows that
\(f\) and \(g\) are equal on every fiber; hence the subprescheme
\(\operatorname{Ker}(f,g)\) of \(X\) is set-theoretically equal to
\(X\), and is therefore closed. Since it contains \(U\), it is equal
to \(X\), that is, \(f=g\). Case \SGARefLink{sga3:xviii:prop:1.4:item:ii}{(ii)} is proved as in
\textit{loc.\ cit.}\footnote{Editor's note: Give details on this
point\ldots}

\SGARefTarget{sga3:xviii:def:1.5}\paragraph{\SGARefLink{sga3:xviii:def:1.5}{Definition 1.5}.}
\begingroup
\renewcommand{\thefootnote}{*}
\footnote{See EGA IV\textsubscript{4}, 20.5, where the term
``pseudo-morphism from \(X\) to \(Y\) relative to \(S\)'' is used in
order not to conflict with EGA I, 7.12.}
\addtocounter{footnote}{-1}
\endgroup
\textit{Let \(X\) be a prescheme over \(S\), and let \(Y\) be a
presheaf over \(S\). A \emph{rational map}
\(f:X\dashrightarrow Y\) over \(S\) is an equivalence class of
morphisms \(f:U\to Y\) over \(S\), where \(U\) is an open subset of
\(X\) that is schematically dense relative to \(S\). One sets
\[
  (f:U\to Y)\sim(f':U'\to Y)
\]
if and only if there exists an open subset
\(U''\subseteq U'\cap U\), schematically dense relative to \(S\),
such that \(f=f'\) on \(U''\).}

This definition is made in such a way that, for every base change
\(S'\to S\), the rational map \(f\times_S S'\) may be defined in the
obvious way.

If \(U\) is an open subset of \(X\), we say that the rational map
\(f\) is \emph{defined on \(U\)} if there exists a morphism
representing \(f\) whose set of definition contains \(U\).

\SGARefTarget{sga3:xviii:def:1.5.1}\paragraph{\SGARefLink{sga3:xviii:def:1.5.1}{Definition 1.5.1}.}
\footnote{Editor's note: The number \SGARefLink{sga3:xviii:def:1.5.1}{1.5.1} has been added.}
\textit{If \(Y\) satisfies one of conditions \SGARefLink{sga3:xviii:prop:1.4:item:i}{(i)}, \SGARefLink{sga3:xviii:prop:1.4:item:ii}{(ii)} of \SGARefLink{sga3:xviii:prop:1.4}{1.4}, it is
clear that there exists a largest open subset \(U\) of \(X\) on
which \(f\) is defined, and this \(U\) is schematically dense
relative to \(S\). It is called the \emph{domain of definition of
\(f\) over \(S\)}, and is denoted \(\Dom(f)\).}

This notion does not in general commute with base change, but one
has:

\SGARefTarget{sga3:xviii:prop:1.6}\paragraph{\SGARefLink{sga3:xviii:prop:1.6}{Proposition 1.6}.}
\textit{Let \(X\) be an \(S\)-prescheme and \(Y\) an \(S\)-functor
satisfying one of hypotheses \SGARefLink{sga3:xviii:prop:1.4:item:i}{(i)}, \SGARefLink{sga3:xviii:prop:1.4:item:ii}{(ii)} of \SGARefLink{sga3:xviii:prop:1.4}{1.4}. Let
\(f:X\dashrightarrow Y\) be a rational map over \(S\), let
\(S'\to S\) be a flat morphism locally of finite presentation, and
put \(f'=f\times_S S'\). Then}
\[
  \Dom(f')=\Dom(f)\times_S S'.
\]

\textit{Proof.}
Put \(U=\Dom(f)\). It is clear that \(V'=\Dom(f')\) contains
\(U\times_S S'\). Let \(V\) be the image of \(V'\); it is an open
subset of \(X\), because \(X'\to X\) is open. We must prove that
\(V=U\), that is, we must find a morphism \(V\to Y\) representing
\(f\). Put \(S''=S'\times_S S'\), and
\[
  X''=X\times_S S'',\qquad
  U''=U\times_S S'',\qquad
  V''=V'\times_V V'.
\]
Then \(U''\) is schematically dense in \(X''\) relative to \(S''\);
hence \(U''\) is schematically dense in \(V''\) relative to
\(S''\), since \(U''\subseteq V''\subseteq X''\). The restriction
of \(f':V'\to Y'\) to \(U'\) is obtained from \(f:U\to Y\) by base
change. The two morphisms \(V''\to Y\) obtained from \(f'\) by the
two projections agree on \(U''\), and hence agree.\begingroup
\renewcommand{\thefootnote}{\ensuremath{\dagger}}\footnote{The
French source prints \(S\), rather than \(S''\), in the preceding
density assertion and the undefined symbol \(g\), rather than
\(f'\), in this descent assertion. The base-change construction
determines the readings used here.}
\addtocounter{footnote}{-1}
\endgroup Since \(V'\to V\) is flat and
locally of finite presentation, it is an fppf covering (\SGARefLink{sga3:iv:subsection:6:3}{Exposé IV,
6.3}), and the morphism \(V\to Y\) is obtained by descent.

We shall frequently use the following triviality:

\SGARefTarget{sga3:xviii:prop:1.7}\paragraph{\SGARefLink{sga3:xviii:prop:1.7}{Proposition 1.7}.}
\textit{Let \(X/S\) be faithfully flat and locally of finite
presentation, and let \(U\subseteq X\) be a relatively schematically
dense open subset. Then there exists an fppf covering base change
\(S'\to S\), and a section \(x\in X(S')\) contained in \(U(S')\).}

Indeed, \(U/S\) is an fppf covering. Take \(S'=U\), and as the section
take the \(S\)-graph of the inclusion \(U\hookrightarrow X\).
